\documentclass[12pt]{article}
\usepackage[utf8]{inputenc}
\usepackage[T1]{fontenc}
\usepackage{hyperref}
\title{Lalu, Eastern Language Resources for K-12 Curriculum Development}
\author{Compiled by Claude Sonnet 5.5 (claude-sonnet-5-5) for Andr\'as Kornai}
\date{}
\begin{document}
\maketitle

\section{Introduction}
Eastern Lalu (ISO 639-3 yit; Glottocode east2696) is a Sino-Tibetan (Loloish) language of Yunnan, China. Searches found almost no dedicated material for it. Sources also disagree on how the Lalo/Lalu ISO codes map to varieties, so check the code mapping before relying on any one source. This report lists only what the searches returned.

\section{Primary Educational Resources}
No primers, textbooks or school programmes for Eastern Lalu were found.

\section{Linguistic and Reference Materials}
\begin{itemize}
\item Wikipedia, Lalo language (general overview of the Lalo cluster and its ISO codes): \url{https://en.wikipedia.org/wiki/Lalo_language}
\item Endangered Languages Project entry for Central Lalo (a related variety, not Eastern Lalu): \url{https://endangeredlanguages.com/lang/5129}
\item Max Planck EVA page on Western Lalu numerals (a related variety, not Eastern Lalu): \url{https://lingweb.eva.mpg.de/channumerals/Lalu-Western.htm}
\end{itemize}
No grammar or dictionary of Eastern Lalu was found.

\section{Religious and Community Resources}
\begin{itemize}
\item Asia Harvest people-group profile, Lalu (Xinping): \url{https://people-groups.asiaharvest.org/people-group-profiles/yunnan/lalu-xinping.pdf}
\end{itemize}

\section{Digital and Audio Resources}
\begin{itemize}
\item Joshua Project language page for yit: \url{https://www.joshuaproject.net/index.php/languages/yit}
\end{itemize}
This page is a directory entry. I did not confirm that it hosts any audio.

\section{Curriculum Development Context}
The reference description says the language is thought to be used as a first language by adults only, and it is not known to be taught in schools. Curriculum work would have to start from fieldwork and documentation, since no teaching materials were found.
\end{document}